% !TeX root = ../../paper.tex

\subsection{Limitations}
\label{sec:discussion-limitations}

Several constraints qualify our results. The sample is small: seven participants, of whom five completed both conditions, so the paired analysis rests on five pairs, and the exact Wilcoxon test cannot reach significance at that size. The inferential statistics orient; they prove nothing. We collected no demographic data beyond an observed age range and a gender count. Because the two single-condition participants played the haptic version, the descriptive condition means also draw on different groups. We measured presence with a single questionnaire; without an emotional measure such as the Self-Assessment Manikin or physiological data such as heart rate variability, the emotional side of the question (does the touch make it scarier, beyond making it more present?) rests on anecdote. The IPQ is a blunt instrument for a manipulation aimed at one dimension above all~\cite{presence-questionnaires}.

The within-subject design carries a specific cost in a horror context: a participant who plays the haptic condition second knows the scares by then, and one who plays it first may brace for touches that the baseline never delivers. Counterbalancing distributes this effect across conditions without removing it. The pipe challenge adds a human factor, since a researcher operated the water spray by hand and its timing and intensity varied between sessions. The study procedure was informal in ways a full research project would not permit: participants gave verbal consent with no written forms, and the study room was a regular work room, so ambient noise and the visible presence of the experimenters may have pulled participants out of the experience in ways we could not measure. \todoin{Add anything you noticed during the sessions---tracking hiccups, a proxy that repeatedly misaligned, participants peeking under the headset.}
